\chapter{Background}\label{chapter:background}

This chapter explains the main concepts used in this thesis. It first
introduces ROS~2 and its communication mechanisms. It then explains runtime
monitoring and its basic terminology, describes how ROS~2 applications can be
observed, and summarizes the main tasks of a monitoring system. Finally, it
introduces feature models as a way to describe variability, and discusses
deployment and transport choices.

\section{ROS 2 Applications}

\subsection{Nodes and the ROS Graph}

ROS~2 is a middleware and software platform for building robotic applications.
Despite its name, it is not an operating system such as Linux. It runs on top
of a conventional operating system (usually Ubuntu) and provides communication, discovery,
build tools, and many reusable packages that help software components work
together~\cite{macenski2022ros2}.

A ROS~2 application is usually divided into \emph{nodes}. Each node has a
focused task. For example, one node may publish camera data, another node may
plan a route, and another node may control the motors. The active nodes and
their communication links form the \emph{ROS graph}. Nodes can run on the same
computer or on different computers.

ROS~2 communication uses defined types. A type describes the fields of the
data that nodes exchange, and client libraries such as \texttt{rclcpp} and
\texttt{rclpy} provide the matching C++ and Python APIs. Nodes communicate
through three main mechanisms: topics, services, and actions. The following
subsections explain each mechanism.

\subsection{Topics}

Topics provide publish-subscribe communication. A publisher sends typed
messages to a topic name. Any number of subscribers can receive these messages.
The publisher does not need to know the subscribers directly. Topics are often
used for sensor data, robot state, movement commands, and
diagnostics~\cite{ros2interfaces,ros2topics}. Figure~\ref{fig:topic-comm}
shows the communication pattern: one publisher sends messages to a topic, and
every subscriber of that topic receives them.

\begin{figure}[htbp]
\centering
\begin{tikzpicture}[
    node distance=0.5cm and 1.8cm,
    ros/.style={draw, rounded corners, minimum width=2.6cm,
        minimum height=1.0cm, align=center},
    chan/.style={draw, minimum width=2.8cm, minimum height=1.0cm,
        align=center, fill=black!8},
    msg/.style={-{Stealth[length=2.5mm]}, thick}]
  \node[chan] (topic) {Topic\\\texttt{/odom}};
  \node[ros, left=of topic] (pub) {Publisher\\node};
  \node[ros, above right=of topic] (sub1) {Subscriber\\node A};
  \node[ros, below right=of topic] (sub2) {Subscriber\\node B};
  \draw[msg] (pub) -- node[above] {\small publish} (topic);
  \draw[msg] (topic.east) -- node[above, sloped] {\small message} (sub1.west);
  \draw[msg] (topic.east) -- node[below, sloped] {\small message} (sub2.west);
\end{tikzpicture}
\caption{Topic communication. A publisher sends typed messages to a named
topic, and all subscribers of that topic receive the messages.}
\label{fig:topic-comm}
\end{figure}

Messages on a topic flow in one direction, from publishers to subscribers.
Several nodes may publish to the same topic, and a node can be a publisher on
some topics and a subscriber on others at the same time.

\subsection{Services}
\label{sec:ros2-services}

Services provide request-response communication. A client sends a request to a
named service, and the server that offers this service returns a response.
Services are normally used for short operations that need a clear result, such
as resetting a pose or switching a component on or
off~\cite{ros2interfaces}. Figure~\ref{fig:service-comm} shows the pattern.

\begin{figure}[htbp]
\centering
\begin{tikzpicture}[
    ros/.style={draw, rounded corners, minimum width=3.0cm,
        minimum height=1.4cm, align=center},
    msg/.style={-{Stealth[length=2.5mm]}, thick}]
  \node[ros] (client) {Service client\\node};
  \node[ros, right=5.2cm of client] (server) {Service server\\node};
  \draw[msg] ([yshift=3.5mm]client.east) --
      node[above] {\small request} ([yshift=3.5mm]server.west);
  \draw[msg] ([yshift=-3.5mm]server.west) --
      node[below] {\small response} ([yshift=-3.5mm]client.east);
\end{tikzpicture}
\caption{Service communication. A client sends a request to a named service,
and the server returns a response to that client.}
\label{fig:service-comm}
\end{figure}

In contrast to topics, service communication is a one-to-one exchange. The
request is delivered only to the server, and the response is delivered only to
the client that sent the request. This property matters later in this chapter,
because it limits how services can be observed.

\subsection{Actions}

Actions are used for tasks that take time to complete. An action client sends
a goal to an action server. While the goal is running, the server sends
feedback and status updates. The client can cancel the goal, and the server
later returns a result. Navigation and robot manipulation are common
examples~\cite{ros2interfaces,ros2actions}. Figure~\ref{fig:action-comm} shows
the pattern.

\begin{figure}[htbp]
\centering
\begin{tikzpicture}[
    ros/.style={draw, rounded corners, minimum width=3.2cm,
        minimum height=2.9cm, align=center},
    msg/.style={-{Stealth[length=2.5mm]}, thick},
    stream/.style={-{Stealth[length=2.5mm]}, thick, dashed}]
  \node[ros] (client) {Action client\\node};
  \node[ros, right=6.2cm of client] (server) {Action server\\node};
  \draw[msg] ([yshift=10.5mm]client.east) --
      node[above] {\small goal request, cancel (services)}
      ([yshift=10.5mm]server.west);
  \draw[stream] ([yshift=3.5mm]server.west) --
      node[above] {\small feedback (topic)} ([yshift=3.5mm]client.east);
  \draw[stream] ([yshift=-3.5mm]server.west) --
      node[above] {\small status (topic)} ([yshift=-3.5mm]client.east);
  \draw[msg] ([yshift=-10.5mm]server.west) --
      node[below] {\small result (service)} ([yshift=-10.5mm]client.east);
\end{tikzpicture}
\caption{Action communication. Goal, cancellation, and result exchanges use
services, while feedback and status updates use topics (dashed arrows).}
\label{fig:action-comm}
\end{figure}

Internally, an action is built from hidden topics and services. Goal, result,
and cancellation communication use services. Feedback and status updates use
topics~\cite{ros2actions}. An action therefore combines the one-to-one
request-response pattern with continuous one-way data streams.

\subsection{Discovery, Middleware, and Quality of Service}
\label{sec:ros2-discovery}

ROS~2 commonly uses Data Distribution Service (DDS) middleware for
communication. DDS provides distributed discovery, so ROS~2 does not need the
central master used in ROS~1~\cite{macenski2022ros2}. Nodes can discover the
names and types of visible topics and services at runtime. Discovery takes
time, however, and the ROS graph can change while the application is running,
so the set of visible names and types is not a fixed snapshot.

ROS~2 also provides Quality of Service (QoS) settings. These settings control
reliability, stored history, durability, deadlines, and other communication
behaviour~\cite{ros2qos}. Two endpoints can only communicate when their QoS
settings are compatible.

\section{Runtime Monitoring and Runtime Verification}
\label{sec:runtime-monitoring}

\subsection{Main Terms}

Runtime monitoring means observing a system while it runs. It can be used for
many purposes, including operational diagnosis, error detection, security, and
requirements checking~\cite{rabiser2019domain}.

Runtime verification is a related technique. It compares an observed execution
with a specification~\cite{leucker2009brief}. It can show whether the observed
behaviour follows or breaks a property; it does not normally prove that every
possible future execution will be correct. In the literature, the two terms
are very often used as synonyms. In this thesis, runtime monitoring is the
broader activity of observing and processing runtime data, and property
checking is one of the tasks it can support.

This thesis uses the following terms:

\begin{description}
    \item[System under monitoring] The ROS~2 application being observed.

    \item[Observation or event] One captured runtime event, such as a topic
    message or service request.

    \item[Trace] A sequence of observations used for later checking.

    \item[Property] A rule or expected condition.

    \item[Monitor] A component that observes the system or checks properties.

    \item[Verdict] The result of a property check.
\end{description}

Figure~\ref{fig:monitoring-terms} shows how these terms relate. The monitor
captures observations from the system under monitoring. Based on a property,
it evaluates the observed behaviour and reports a verdict to a user or to
another component.

\begin{figure}[htbp]
\centering
\begin{tikzpicture}[
    box/.style={draw, rounded corners, minimum width=3.0cm,
        minimum height=1.2cm, align=center},
    msg/.style={-{Stealth[length=2.5mm]}, thick}]
  \node[box] (system) {System under\\monitoring};
  \node[box, right=2.6cm of system] (monitor) {Monitor};
  \node[box, right=2.6cm of monitor] (user) {User or\\component};
  \node[box, dashed, above=1.0cm of monitor] (property) {Property};
  \draw[msg] (system) -- node[above] {\small observations}
      node[below] {\small (trace)} (monitor);
  \draw[msg] (monitor) -- node[above] {\small verdict} (user);
  \draw[msg] (property) -- (monitor);
\end{tikzpicture}
\caption{Basic runtime monitoring terminology. A monitor captures observations
from the system under monitoring, checks them against a property, and reports
a verdict.}
\label{fig:monitoring-terms}
\end{figure}

\subsection{Online and Offline Monitoring}

Online monitoring processes data while the system is running. It can report
problems quickly and can support immediate feedback. The design question is
where this monitoring logic should run and which observations need to be
processed online.

Offline monitoring checks recorded data after the run has finished. It cannot
react immediately, but the same data can be replayed many times and checked
with checkers that are easier to run outside the robot application.

\subsection{Passive and Active Monitoring}

A passive monitor observes and reports. It does not change the behaviour of the
application. An active monitor can respond to a problem, for example by
blocking a message or starting a recovery action.

\subsection{Centralized and Distributed Monitoring}

In a centralized setup, observations are sent to one main checking component.
Keeping the checking state in one place simplifies the design, although the
central component can become a bottleneck or a single point of failure.

In a distributed setup, observation or checking work runs in several places.
This can place monitoring logic closer to the data source or closer to a
verifier, but it introduces questions about message ordering, clocks, network
boundaries, and failures~\cite{francalanza2018distributed}. Concrete
deployment diagrams for the setups used in this thesis are shown in
Chapter~\ref{chapter:design}.

\section{Observing ROS 2 Applications}
\label{sec:observing-ros2}

A monitor for a ROS~2 application first has to obtain observations. The three
communication mechanisms from the previous section can each be observed, but
not in the same way. This section explains the observation mechanism for each
of them, and the role that types and discovery play for a monitor.

\subsection{Observing Topics}

A monitor can observe a topic by creating an extra subscription, as any other
node would. This is a simple and passive method because the original
communication path does not need to change. Common ROS~2 tools such as
\texttt{ros2 topic echo} and \texttt{ros2 bag record} work in this way.

Topic observation still requires several choices. Some topics publish large
messages or publish at a high rate, while a monitor may only need a few fields
or a subset of the messages. The monitoring subscription must also use a QoS
setting that is compatible with the publisher; otherwise it may receive no
messages or lose data under high load. Finally, messages from different
publishers may arrive in a different order from the order in which they were
sent.

\subsection{Observing Services}
\label{sec:observing-services}

Services cannot be observed in the same way as topics. As explained in
Section~\ref{sec:ros2-services}, a service call is a one-to-one exchange: the request travels
from one client to the server, and the response travels back to that client
only. A third node that wants to watch the service cannot ``listen in'' on
this exchange. If it creates its own client for the same service, it can only
make its own calls; it never sees the requests and responses of other clients.

ROS~2 solves this with \emph{service introspection}. When introspection is
enabled on the client or the server, that node publishes an event message on
an additional \emph{service event topic} for every step of a call. Each event
carries metadata (such as the event type, a client identifier, a sequence
number, and a timestamp) and can also carry the request or response
payload~\cite{ros2serviceintrospection}. Because the events are published on a
normal topic, a monitor can observe them with a plain subscription, as shown
in Figure~\ref{fig:service-introspection}.

\begin{figure}[htbp]
\centering
\begin{tikzpicture}[
    ros/.style={draw, rounded corners, minimum width=2.9cm,
        minimum height=1.2cm, align=center},
    chan/.style={draw, minimum width=5.2cm, minimum height=1.1cm,
        align=center, fill=black!8},
    msg/.style={-{Stealth[length=2.5mm]}, thick},
    event/.style={-{Stealth[length=2.5mm]}, thick, dashed}]
  \node[ros] (client) {Service client};
  \node[ros, right=5.6cm of client] (server) {Service server};
  \node[chan] (evtopic) at ($(client)!0.5!(server) + (0,-2.6)$)
      {Service event topic\\\texttt{<service\_name>/\_service\_event}};
  \node[ros, below=1.1cm of evtopic] (monitor) {Monitor};
  \draw[msg] ([yshift=3.5mm]client.east) --
      node[above] {\small request} ([yshift=3.5mm]server.west);
  \draw[msg] ([yshift=-3.5mm]server.west) --
      node[below] {\small response} ([yshift=-3.5mm]client.east);
  \draw[event] (client.south) |- node[below left, pos=0.75]
      {\small event messages} (evtopic.west);
  \draw[event] (server.south) |- node[below right, pos=0.75]
      {\small event messages} (evtopic.east);
  \draw[event] (evtopic) -- node[right] {\small subscribe} (monitor);
\end{tikzpicture}
\caption{Observing a service with service introspection. The request and
response are exchanged only between the client and the server (solid arrows).
With introspection enabled, the client and server additionally publish event
messages on a service event topic, which a monitor can observe through a
normal subscription (dashed arrows).}
\label{fig:service-introspection}
\end{figure}

This mechanism keeps the observation passive, but it has two consequences for
a monitor. First, introspection must be enabled by the application; without
it, the event topic carries no data. Second, a request and its response arrive
as separate event messages, so the monitor must use the event metadata, such
as the client identifier and sequence number, to connect a response to the
correct request.

\subsection{Observing Actions}

Actions combine the previous two cases, because they are built from hidden
topics and services. The feedback and status streams are topics, so a monitor
can subscribe to them directly. The goal, result, and cancellation exchanges
are services, so observing them requires service introspection on the action
server or client.

Action observation adds one practical concern: the hidden topic and service
names are internal ROS~2 details. A monitor should present the captured events
under the logical action name, such as \texttt{/navigate\_to\_pose}, instead
of forcing later components to understand all hidden names.

\subsection{Types and Discovery for Monitors}

Type information is important for every observation mechanism above. A monitor
needs the correct type before it can read and understand captured data.
Discovery can help a monitor find the type of a topic or service
automatically, but as explained in Section~2.1.5, discovery takes time and the
graph can change at runtime, so a monitor should not assume that every source
is already visible at startup.

The ROS graph also shapes what monitoring looks like in ROS~2. A monitoring
node can often join the graph without changing the original application, which
makes passive observation natural. At the same time, the behaviour of the full
robot is spread over many nodes, so a monitor must select useful interfaces
and keep enough context to connect related events.

\section{Main Monitoring Tasks}

A monitoring system usually performs several different
tasks~\cite{rabiser2019domain}. Keeping these tasks separate makes the
infrastructure easier to configure and extend.

\subsection{Collecting Data}

The first task is to obtain data from the running system. In ROS~2, this can
mean subscribing to topics, reading service events, observing action feedback,
using lower-level tracing tools, or even instrumenting the application code
(Section~\ref{sec:observing-ros2}).

Each method has different limits. Normal subscriptions are easy to add, but
they only show what the monitor receives. Service introspection needs support
from the client or server. Lower-level tracing can show detailed execution
information but may not include the application meaning of a message.
Instrumenting the code gives the most direct access, but it changes the
application itself.

\subsection{Processing Data}

Raw data is not always ready for checking. A monitor may need to select fields,
rename values, add timestamps, or connect events to a session. It may also need
to reduce the amount of data.

\subsection{Routing and Storage}

After processing, data must be sent to the correct destination. One
observation may be written to a file, sent over a network, and passed to a
checking component at the same time. Adding or changing such destinations
should not require changes to the components that collect the data.

Storage also supports later analysis. A stable data format makes it easier to
replay a run or to compare results from different checking tools.

\subsection{Checking and Verdicts}

Checking decides whether observed data follows an expected property. Different
tools need different input formats. A simple threshold checker may need one
number, while a temporal-logic checker may need a sequence of timestamped
values.

\subsection{Using the Results}

A verdict is only useful when it reaches the right user or component. It may be
printed during development, stored in a file, sent to a dashboard, or passed to
another service.

\section{Variability and Feature Models}
\label{sec:feature-models}

Systems in the same domain usually share a common core but differ in many
details. Building one configurable infrastructure for such a domain only works
when these common and variable parts are identified first: the common parts
can be implemented once, while the variable parts must become configuration
choices. This requires a systematic way to describe what can vary and which
combinations are allowed.

Feature-Oriented Domain Analysis (FODA) provides such a
description~\cite{kang1990foda}. In FODA, the capabilities of the systems in a
domain are described as \emph{features}, organized in a tree called a
\emph{feature model} or feature diagram. The tree shows which features are
always present, which are optional, and which exclude each other, and it can
be extended with cross-tree constraints between features in different
branches~\cite{benavides2010automated}.

Figure~\ref{fig:feature-model-example} shows a small example feature model of
a mobile phone product line, a common illustration in the
literature~\cite{benavides2010automated}. The notation is read as follows:

\begin{itemize}
    \item A \textbf{mandatory} feature (filled circle) is part of every
    product in which its parent appears. Every phone supports calls and has a
    screen.

    \item An \textbf{optional} feature (empty circle) may be included or left
    out. A phone may or may not have GPS or media functions.

    \item An \textbf{alternative} group (empty arc) allows exactly one child
    to be chosen. A screen is basic, colour, or high resolution, but only one
    of these.

    \item An \textbf{or} group (filled arc) allows one or more children to be
    chosen. A phone with media functions has a camera, an MP3 player, or both.

    \item A \textbf{cross-tree constraint} (dashed arrow) restricts
    combinations across branches. A camera requires a high-resolution screen.
    Constraints can also exclude combinations, for example when two features
    cannot be used together.
\end{itemize}

\begin{figure}[htbp]
\centering
\begin{tikzpicture}[
    feat/.style={draw, rounded corners=1pt, minimum height=0.72cm,
        inner xsep=6pt, align=center, font=\small},
    fedge/.style={thick},
    cross/.style={dashed, -{Stealth[length=2.2mm]}}]
  \node[feat] (phone)  at (0,0)        {Mobile phone};
  \node[feat] (calls)  at (-5.2,-1.9)  {Calls};
  \node[feat] (gps)    at (-3.4,-1.9)  {GPS};
  \node[feat] (screen) at (-0.9,-1.9)  {Screen};
  \node[feat] (media)  at (3.8,-1.9)   {Media};
  \node[feat] (basic)  at (-3.0,-3.8)  {Basic};
  \node[feat] (colour) at (-1.1,-3.8)  {Colour};
  \node[feat] (high)   at (1.1,-3.8)   {High res.};
  \node[feat] (camera) at (3.0,-3.8)   {Camera};
  \node[feat] (mp3)    at (4.9,-3.8)   {MP3};
  \draw[fedge] (phone.south) -- (calls.north);
  \draw[fedge] (phone.south) -- (gps.north);
  \draw[fedge] (phone.south) -- (screen.north);
  \draw[fedge] (phone.south) -- (media.north);
  \draw[fedge] (screen.south) -- (basic.north);
  \draw[fedge] (screen.south) -- (colour.north);
  \draw[fedge] (screen.south) -- (high.north);
  \draw[fedge] (media.south) -- (camera.north);
  \draw[fedge] (media.south) -- (mp3.north);
  % mandatory / optional markers
  \fill (calls.north) circle (2.2pt);
  \fill (screen.north) circle (2.2pt);
  \draw[fill=white] (gps.north) circle (2.2pt);
  \draw[fill=white] (media.north) circle (2.2pt);
  % alternative (xor) group: empty arc
  \draw ($(screen.south)!7mm!(basic.north)$) to[bend right=32]
        ($(screen.south)!7mm!(high.north)$);
  % or group: filled arc
  \fill ($(media.south)!6.5mm!(camera.north)$) -- (media.south) --
        ($(media.south)!6.5mm!(mp3.north)$) to[bend right=40] cycle;
  % cross-tree constraint
  \draw[cross] (camera.south) .. controls +(0,-0.9) and +(0,-0.9) ..
      node[below, font=\small\itshape] {requires} (high.south);
\end{tikzpicture}
\caption{Example feature model of a mobile phone product line (adapted
from~\cite{benavides2010automated}). Filled circles mark mandatory features,
empty circles mark optional features, an empty arc marks an alternative
(exactly one) group, a filled arc marks an or (one or more) group, and the
dashed arrow is a cross-tree constraint.}
\label{fig:feature-model-example}
\end{figure}

A concrete product is a selection of features that respects all these rules.
In this thesis, the same notation is used in
Chapter~\ref{chapter:design} to describe the commonalities and variabilities
of ROS~2 monitoring solutions.

\section{Deployment and Transport}

\subsection{Integrated Monitoring}

In the integrated setup, collection, processing, checking, and output run
together. This setup is simple and does not need a network connection between
the monitor and verifier. It is useful for local development and smaller
systems.

The price of this simplicity is tighter coupling: collection and checking share
the same runtime context, so checker-specific failures or incompatible
extension behaviour can affect the full monitoring process.

\subsection{Split Monitoring}

In the split setup, a robot-side process collects and processes observations.
A separate verifier receives the observations and performs checking. This
keeps ROS~2-specific collection close to the robot while placing
checker-specific logic in a separate process.

This setup introduces a network boundary. Data may be delayed, lost, or
received in a different order. The verifier also needs enough information in
each observation to understand its source and session. Deployment diagrams
for both setups are shown in Chapter~\ref{chapter:design}.

\subsection{Transport Protocols}

When monitoring components run in different processes or on different
machines, as in the split setup, they need a transport mechanism to exchange
data. Many protocols can fill this role. The components could communicate
through the DDS middleware that ROS~2 itself uses, through plain TCP or UDP
sockets, through web protocols such as HTTP or WebSocket, or through message
brokers based on protocols such as AMQP, Kafka, or MQTT. These options differ
in how they handle addressing, delivery guarantees, and decoupling between
sender and receiver.

MQTT is a publish-subscribe messaging protocol that is commonly used in
robotics and Internet-of-Things systems, because it is simple, lightweight,
and has wide tool support. Publishers send messages to named topics through a
broker, and subscribers receive messages from the topics they select. MQTT~5.0
is an OASIS standard~\cite{mqtt5spec}. Because the broker decouples publishers
from subscribers, neither side needs to know the network address of the other.
This thesis uses MQTT as the transport in the implementation of the case
studies; the transport itself remains a configurable choice in the
infrastructure.

In the split monitoring deployment, the robot-side monitor publishes
serialized observations to a configured MQTT topic or topic prefix. The
verifier subscribes to that topic filter and receives the data without needing
a ROS~2 dependency or a direct connection to the robot process. Verdicts can
also be published to a separate MQTT topic, which keeps runtime data and
checking results distinct.

MQTT also gives a configurable delivery trade-off through Quality of Service
(QoS) levels. QoS~0 sends a message at most once and is useful when occasional
loss is acceptable. QoS~1 retries until the broker acknowledges the message,
which can duplicate messages and therefore requires identifiers to make
duplicates recognizable. QoS~2 adds a stronger exactly-once exchange at the
expense of a more involved protocol exchange. The infrastructure keeps the
session and record identifier inside each observation so that the transport
does not have to carry the semantic identity of an observation.
